\documentclass[11pt, a4paper, twoside]{article}

% --- Packages ---
\usepackage[utf8]{inputenc}
\usepackage{amsmath, amssymb, amsthm}
\usepackage{geometry}
\usepackage{tcolorbox}
\usepackage{tikz-cd}
\usepackage{fancyhdr}
\usepackage{booktabs} % For professional tables

% --- Page Layout & Styling ---
\geometry{
	top=1in, 
	bottom=1in, 
	left=1.25in, 
	right=1.25in
}

\pagestyle{fancy}
\fancyhf{}
\fancyhead[LE,RO]{\thepage}
\fancyhead[LO]{\small \textit{Lecture Notes: Homological Coding Theory}}
\fancyhead[RE]{\small \textit{The de Rham Perspective}}
\renewcommand{\headrulewidth}{0.4pt}

% Custom boxes
\newtcolorbox{theorybox}[1][]{
	colback=gray!5!white, colframe=gray!75!black,
	fonttitle=\bfseries, title=#1, arc=3pt, boxrule=1pt,
	left=10pt, right=10pt, top=10pt, bottom=10pt
}
\newtcolorbox{examplebox}[1][]{
	colback=blue!5!white, colframe=blue!75!black,
	fonttitle=\bfseries, title=#1, arc=3pt, boxrule=1pt,
	left=10pt, right=10pt, top=10pt, bottom=10pt
}

% --- Macros ---
\newcommand{\F}{\mathbb{F}}
\newcommand{\im}{\text{im}}
\newcommand{\kerop}{\text{ker}}

\begin{document}
	
	\begin{center}
		{\Large \textbf{Geometricizing Codes: The de Rham Perspective}} \\[0.5em]
		{\normalsize \textit{Translating Linear Algebra into Differential Geometry}}
	\end{center}
	\vspace{1em}
	
	\section*{1. The Translation Dictionary}
	
	To view a linear block code as a geometric object, we must translate the discrete finite-field mechanics of coding theory into the continuous differential geometry of manifolds. We map the sequence of vector spaces over $\F_q$ directly onto the de Rham complex of differential forms $\Omega^i(M)$ on a smooth manifold $M$.
	
	\begin{table}[h!]
		\centering
		\renewcommand{\arraystretch}{1.4}
		\begin{tabular}{@{}ll@{}}
			\toprule
			\textbf{de Rham Cohomology} & \textbf{Linear Coding Theory} \\
			\midrule
			\textbf{0-forms} $\Omega^0(M)$ (Smooth functions) & \textbf{Message Space} $\F_q^k$ (Raw information) \\
			\textbf{1-forms} $\Omega^1(M)$ (Vector fields/1-forms) & \textbf{Ambient Space} $\F_q^n$ (All possible words) \\
			\textbf{2-forms} $\Omega^2(M)$ (Area forms) & \textbf{Syndrome Space} $\F_q^{n-k}$ (Parity violations) \\
			\textbf{Exterior Derivative} $d_0: \Omega^0 \to \Omega^1$ & \textbf{Generator Matrix} $G: \F_q^k \to \F_q^n$ \\
			\textbf{Exterior Derivative} $d_1: \Omega^1 \to \Omega^2$ & \textbf{Parity-Check Matrix} $H: \F_q^n \to \F_q^{n-k}$ \\
			\textbf{Nilpotency} $d_1 \circ d_0 = 0$ ($d^2=0$) & \textbf{Orthogonality} $H G^T = \mathbf{0}$ \\
			\bottomrule
		\end{tabular}
	\end{table}
	
	\section*{2. Exactness and Classical Codes}
	
	We align the de Rham complex and the coding complex side-by-side:
	
	\[
	\begin{tikzcd}[row sep=huge, column sep=large]
		0 \arrow[r] & 
		\Omega^0(M) \arrow[r, "d_0"] \arrow[d,leftrightarrow, dashed] & 
		\Omega^1(M) \arrow[r, "d_1"] \arrow[d,leftrightarrow, dashed] & 
		\Omega^2(M) \arrow[d,leftrightarrow, dashed] \\
		0 \arrow[r] & 
		\F_q^k \arrow[r, "G"] & 
		\F_q^n \arrow[r, "H"] & 
		\F_q^{n-k} 
	\end{tikzcd}
	\]
	
	\begin{theorybox}[Classical Codes are Topologically Trivial]
		In a standard classical block code, the sequence is strictly \textbf{exact} at the ambient space: $\im(G) = \kerop(H)$. 
		\begin{itemize}
			\item $\kerop(H)$ represents the \textbf{closed 1-forms} (words with no parity violations; valid codewords).
			\item $\im(G)$ represents the \textbf{exact 1-forms} (words generated from a raw message).
		\end{itemize}
		Because the sequence is exact, every closed form is exact. There are no topological ``holes.'' The cohomology group is trivial: $H^1_{dR} = \kerop(H) / \im(G) = 0$. The valid codewords are exactly the generated messages.
	\end{theorybox}
	
	\section*{3. Topological Codes: Cohomology as Information}
	
	If we deliberately design a coding sequence that is a chain complex ($H \circ G = 0$) but is \textbf{not exact}, we create a Calderbank-Shor-Steane (CSS) Topological Quantum Code. By relaxing exactness, $\im(G) \subsetneq \kerop(H)$, and we create a non-trivial cohomology group:
	$$ H^1 = \frac{\kerop(H)}{\im(G)} \neq 0 $$
	
	\begin{itemize}
		\item \textbf{The Code Space ($\kerop(H)$):} The closed forms. Physically, these are quantum states satisfying all stabilizers.
		\item \textbf{Gauge Transformations ($\im(G)$):} The exact forms. Local, trivial fluctuations that do not change stored data.
		\item \textbf{Logical Information ($H^1$):} The true protected information is stored in the cohomology classes. Two codewords $c_1, c_2 \in \kerop(H)$ encode the exact same logical information if they differ only by an exact form ($c_1 - c_2 \in \im(G)$).
	\end{itemize}
	
	\section*{4. Concrete Construction: The Toric Code}
	
	\begin{examplebox}[Matrix Operators from a 2D Manifold]
		Let $X$ be a triangulation (or a square lattice) of a Torus, which is a 2D manifold. We define our vector spaces over $\F_2$ based on the simplices of the manifold:
		\begin{itemize}
			\item $C_0$: Vector space spanned by the \textbf{Vertices} ($V$).
			\item $C_1$: Vector space spanned by the \textbf{Edges} ($E$).
			\item $C_2$: Vector space spanned by the \textbf{Faces} ($F$).
		\end{itemize}
		We define the coboundary operators $d_0: C_0 \to C_1$ and $d_1: C_1 \to C_2$.
		\begin{align*}
			G &= d_0 \quad (\text{Maps a vertex to its adjacent edges}) \\
			H &= d_1 \quad (\text{Maps an edge to its adjacent faces})
		\end{align*}
		Because the boundary of a boundary is zero, $d_1 \circ d_0 = 0$, guaranteeing $HG^T = \mathbf{0}$. The number of protected logical qubits $k_{logical}$ is given precisely by the Betti number of the manifold:
		$$ k_{logical} = \dim H^1(X, \F_2) $$
		For a Torus, the first Betti number is $b_1 = 2$. Thus, this geometric matrix construction perfectly encodes 2 logical qubits into the global topological ``holes'' of the lattice, protected against any local edge errors (exact forms).
	\end{examplebox}
	
\end{document}